\section{Idea inicial}
\label{sec:Fase0Resumen}

La primera idea fue seguir al pie de la letra el enunciado propuesto:
usar un motor de estados genérico que moviera la lógica del juego, para
lo que habría que diseñar un fichero \emph{JSON} con los datos del
«cartucho»: la imagen de fondo, la lista de los segmentos a usar, la
lista de las imágenes de cada segmento y una tabla con la definición
del motor de estados.

Con esto, el sistema se compondría de un director del juego, que fuera
llamando al intérprete de los estados para definir las transiciones que
pusieran visibles o invisibles a los segmentos, y luego al motor gráfico,
que presentara los visibles y eliminara los invisibles.

El primer problema surgió al pensar que era muy sencillo tomar una
imagen con todos los segmentos, recortarlos manualmente de forma
individual y generar los ficheros de imagen separados. Al comenzar esta
separación aparecieron varios temas:

\begin{itemize}
	\item En el Parachute hay unos 40 segmentos, y en otros juegos
	pueden existir muchos más: un trabajo mucho más laborioso de lo que
	inicialmente parecía.
	\item Al recortar se ve que no pueden ser rectángulos simples: hay
	solapamiento de los recuadros, lo que obliga a retocar algunas
	imágenes después.
	\item Los programas de dibujo aplican suavizado (\emph{anti-aliasing})
	para mejorar el resultado, algo que aquí no interesa: tras recortar
	habría que pasar cada imagen por un proceso que la convierta a blanco
	y negro puros.
\end{itemize}

Por tanto, \underline{deseché esta idea y pensé en un diseñador gráfico
	propio}, que a partir de la imagen en blanco y negro de los segmentos
facilitara todo el proceso.

Para que esto funcionara, las imágenes debían recortarse y escalarse al
tamaño de la pantalla del ESP32, por lo que \underline{era necesario
	añadir una configuración} con ese tamaño en algún lugar; y al pensarlo
vi que también había que definir cómo están unidos los botones al
ESP32. Esa configuración es lo que acabaría siendo el \emph{setup}.

Con estas dos ideas iniciales comencé el desarrollo de la primera fase,
el generador gráfico, pensando que podría añadir luego la tabla de
estados al fichero del cartucho. Pero hay muchas variables a considerar
en el juego y muchos eventos externos a los propios segmentos que
tratar: cuántos se pueden lanzar a la vez en cada iteración, cómo se
definen los que se lanzan inicialmente, las aceleraciones\ldots{} muchos
detalles que hay que definir. Las transiciones de estados, por tanto, no
eran sencillas de definir: aunque se añadan, por ejemplo, 5 variables a
la tabla, con 40 segmentos habría que montar una tabla de $45 \times 45$
entradas, lo que hacía muy compleja su definición.

Por tanto deseché también una máquina de estados pura junto a su tabla
de transiciones, y la cambié por definir cómo se comportan los
segmentos usando una serie de reglas sencillas, que se pueden definir
también en el gráfico de los segmentos. De esta forma, el entorno de
autoría es el propio diseñador gráfico, y no requiere manipular los
ficheros JSON del cartucho: todo se define usando el propio gráfico
como punto de partida.

\clearpage